\scanpage{174}
\mspage{90}
% archive sheet number: 173
\deleted{Ce lemme est un}

% lines 1-4
Démonstration de 4.4.8. Revenons au diagramme (4.4.6.2), dont on déduit (4.4.6.3), nous allons le réécrire de façon simplifiée
\diagram{\begin{tikzcd}[ampersand replacement=\&,column sep=large]
B \arrow[r,"\beta"] \arrow[d,no head,"\varepsilon"] \& B' \arrow[d,no head,"\varepsilon'"] \& \ni y'_0 \arrow[d,mapsto] \\
C \arrow[u,bend left=45,dashed,"\lambda"] \arrow[r,no head,"\gamma"] \& C' \arrow[u,bend left=45,dashed,"\lambda'"] \& \ni z'_0
\end{tikzcd}}
\[\tag*{(4.4.9.1)}\vphantom{B}\]

% lines 5-9
Je dis que $\varepsilon,\varepsilon'$ ont \uline{des adjoints à gauche} $\lambda$ \deleted{$\unsure{\lambda'}$} resp. $\lambda'$, et que l'hom canonique de ``changement de base''
\[
\lambda'\deleted{\unsure{\beta}}\gamma\longrightarrow\beta\lambda\tag*{(4.4.9.2)}
\]
est un \uuline{iso.}

% lines 10-20
En effet, l'existence des adjoints de $\varepsilon,\varepsilon'$, à savoir des foncteurs $\varepsilon_{j!}$ et $\varepsilon'_{j!}$ relatifs à l'inclusion
\[
\varepsilon_j:e\longrightarrow J
\]
et aux catégories de foncteurs $\underline{\operatorname{Hom}}(J,C)$ resp. $\underline{\operatorname{Hom}}(J,\operatorname{Ind}'_{\pi}(C))$, \unsure{résultent|résulte} de \deleted{celle} l'existence des sommes de type $\operatorname{Hom}(j,j')$ ($j'\in\operatorname{Ob}J$) dans $C$ \inserted{(déjà acquise)} \uline{et dans} $\operatorname{Ind}'_{\pi}C$; et du fait que le foncteur $C\longrightarrow\operatorname{Ind}'_{\pi}C$ commute
